\chapter{Heavy Periods}
\index{Heavy Periods}

\section*{Definition and Overview}
Heavy menstrual bleeding, also known as menorrhagia, is menstrual blood loss that is excessive in amount, duration, or both, and that interferes with physical, social, or emotional wellbeing. Around one in five women experiences heavy periods, which can cause anaemia, fatigue, and significant disruption to work, study, and daily life. Heavy bleeding is often due to hormonal imbalance or benign conditions such as fibroids and polyps, and in most cases effective treatment is available. The aim of care is to reduce blood loss, treat any underlying cause, and restore quality of life.

\section*{Epidemiology}
Heavy menstrual bleeding affects approximately one in five women of reproductive age, making it one of the most common gynaecological problems. It accounts for a large proportion of GP and gynaecology visits and is a leading reason for hysterectomy, though fewer hysterectomies are now performed as alternatives have improved. The prevalence is highest in perimenopausal women, and iron-deficiency anaemia is a common consequence. Despite its frequency, many women do not seek help, accepting heavy bleeding as normal and delaying treatment for years.

\section*{Aetiology and Pathophysiology}
\begin{itemize}
    \item \textbf{Hormonal imbalance:} anovulatory cycles, particularly at puberty and perimenopause, lead to unopposed oestrogen and heavy shedding
    \item \textbf{Fibroids:} benign muscle tumours of the uterus that enlarge the cavity and increase bleeding
    \item \textbf{Endometrial polyps:} benign growths that can bleed
    \item \textbf{Adenomyosis:} endometrial tissue within the uterine muscle, causing heavy, painful periods
    \item \textbf{Endometriosis and pelvic infection:} associated with heavy or painful bleeding
    \item \textbf{Bleeding disorders:} von Willebrand disease and platelet disorders in a small proportion of women
    \item \textbf{Thyroid disease and other medical conditions:} can disturb menstrual blood loss
    \item \textbf{Medication:} anticoagulants and some other drugs increase bleeding
\end{itemize}

\section*{Clinical Features}
\begin{itemize}
    \item Soaking through pads or tampons every hour or two
    \item Periods lasting longer than seven days
    \item Large clots and flooding, with blood pooling or leaking
    \item The need to change protection overnight
    \item Periods that interfere with work, school, social life, or intimacy
    \item Symptoms of anaemia: fatigue, weakness, breathlessness, dizziness, and pale skin
    \item Menstrual pain may accompany heavy bleeding
\end{itemize}

\section*{Differential Diagnosis}
\begin{itemize}
    \item Fibroids and polyps are the most common structural causes
    \item Adenomyosis and endometriosis cause heavy, painful periods
    \item Hormonal dysfunction from perimenopause, PCOS, or thyroid disease
    \item Bleeding disorders and anticoagulant use
    \item Early pregnancy complications, including miscarriage, must be excluded in women of reproductive age
    \item Rarely, endometrial hyperplasia or cancer, particularly in older women or those with risk factors
\end{itemize}

\section*{When to Seek Medical Care}
Any woman with periods heavy enough to disrupt her life, cause anaemia, or last longer than seven days should see a doctor. New heavy bleeding in a woman over 40, bleeding between periods, or bleeding after sex also warrants review.

\textbf{Contact your local emergency services for:}
\begin{itemize}
    \item Bleeding so heavy that pads or tampons are soaked within an hour, especially with faintness or dizziness
    \item Severe pelvic pain with bleeding
    \item Heavy bleeding with signs of shock: pale, sweaty, rapid pulse, or collapse
    \item Heavy bleeding during pregnancy
    \item Passing large clots with severe pain
\end{itemize}

\section*{Diagnosis and Investigations}
\begin{itemize}
    \item \textbf{History:} pattern and volume of bleeding, pain, impact on life, and bleeding risk factors
    \item \textbf{Examination:} abdominal and pelvic examination to detect fibroids, tenderness, or masses
    \item \textbf{Blood tests:} full blood count for anaemia, iron studies, and thyroid function
    \item \textbf{Ultrasound:} transvaginal and transabdominal scanning identifies fibroids, polyps, and thickened endometrium
    \item \textbf{Endometrial biopsy or hysteroscopy:} sampling the uterine lining, especially in women over 40 or with risk factors
    \item \textbf{Coagulation testing:} in women with bleeding from multiple sites or a family history of bleeding disorders
    \item \textbf{Pregnancy test:} always performed where relevant
\end{itemize}

\section*{Management}
\begin{itemize}
    \item \textbf{Medical treatment:} tranexamic acid reduces bleeding, and NSAIDs such as mefenamic acid reduce blood loss and pain
    \item \textbf{Hormonal options:} the combined oral contraceptive pill, the levonorgestrel-releasing intrauterine system (Mirena), and progestogens
    \item \textbf{Iron replacement:} corrects and prevents anaemia
    \item \textbf{Minimally invasive surgery:} removal of polyps, myomectomy for fibroids, and endometrial ablation
    \item \textbf{Surgery:} hysterectomy for women who have completed their family and have not responded to other treatment
    \item \textbf{Treatment of underlying causes:} thyroid disease, bleeding disorders, and infections
    \item \textbf{Individualised care:} choice depends on the cause, fertility wishes, age, and preferences
\end{itemize}

\section*{Complications}
\begin{itemize}
    \item Iron-deficiency anaemia with fatigue, breathlessness, and reduced work capacity
    \item Chronic pain and reduced quality of life
    \item Interference with work, education, relationships, and physical activity
    \item Missed underlying conditions, including endometrial cancer in a small minority
    \item Anaemia in pregnancy if untreated
    \item Treatment-related effects, such as surgical complications and hormonal side effects
\end{itemize}

\section*{Prognosis and Outcomes}
Heavy menstrual bleeding is very treatable. Most women achieve a substantial reduction in bleeding with medical therapy, and the hormonal intrauterine system is highly effective, reducing blood loss by up to 90\%. For women who do not respond to medical treatment, surgical options including ablation and hysterectomy are effective, and most women report significant improvement in quality of life. The underlying cause determines the long-term picture: hormonal causes respond well to treatment, while fibroids and adenomyosis may need ongoing management. Untreated, heavy periods cause progressive anaemia and disability, so seeking help early is important.

\section*{Prevention and Self-Care}
\begin{itemize}
    \item Keep a menstrual diary of flow, duration, and symptoms to share with your doctor
    \item Track pad and tampon use to assess the true volume of bleeding
    \item Maintain adequate iron intake through diet and supplements where needed
    \item Do not accept heavy periods as normal; seek medical review
    \item Take iron with vitamin C to improve absorption
    \item Attend follow-up to monitor anaemia and treatment response
    \item Discuss contraception and fertility plans with your doctor when choosing treatment
\end{itemize}

\section*{Sources and Further Reading}
The following reputable sources provide further information for healthcare students and patients seeking reliable, current advice about heavy periods.

\begin{itemize}
    \item Jean Hailes for Women's Health. \textit{Heavy periods: causes and treatment.} https://www.jeanhailes.org.au/
    \item Healthdirect. \textit{Heavy periods (menorrhagia).} https://www.healthdirect.gov.au/heavy-periods
    \item The Royal Women's Hospital. \textit{Heavy menstrual bleeding.} https://www.thewomens.org.au/
    \item NHS. \textit{Heavy periods: overview and treatment.} https://www.nhs.uk/conditions/heavy-periods/
    \item Mayo Clinic. \textit{Menorrhagia: symptoms and causes.} https://www.mayoclinic.org/
    \item MSD Manual. \textit{Heavy menstrual bleeding.} https://www.msdmanuals.com/
\end{itemize}
